\section{准分布与赝分布：异同之比较}

在最近的论文~\cite{Radyushkin:2017cyf} 中，对 LaMET 所提出的欧氏矩阵元给出了另一种诠释，
为从同样的格点计算中提取物理部分子分布提供了另一条途径。出发点是所谓的 Ioffe 时间分布
${\cal M}(\nu,\xi^2)$，它定义为双线性夸克算符在时空间隔 $\xi^\mu$ 任意分隔、处于四动量为
$p^\mu$ 的强子态中的矩阵元。它是 Ioffe 时间 $\nu= -p\cdot \xi$ 与时空不变量 $-\xi^2$ 的
函数。实践中，这样的洛伦兹不变函数可以在欧氏格点上由两夸克空间关联函数
\begin{equation}\label{eq:covitd}
h(zp_z,z^2,\Lambda) =\langle p|\overline{\psi}(z)\Gamma L(z,0)\psi(0)|p\rangle \ ,
\end{equation}
得到，其中 $\Gamma$ 是狄拉克矩阵，$L(z,0)$ 是直规范连接，$\Lambda$ 是重正化标度或格点截断标度。
于是，LaMET 准分布计算中的矩阵元 $h(zp_z,z^2,\Lambda)$ 可视为协变 Ioffe 时间分布
$h(\nu,\xi^2,\Lambda)$，其中 $\nu = zp_z$、类空间隔 $-\xi^2 = z^2$。

Radyushkin 建议在一个新参考系中解释 Ioffe 时间分布：$\xi^\mu = (\xi^+=0, \xi^-/\gamma, \xi_\perp)$，
其中我们使用了光锥变量 $\xi^\mu = (\xi^+,\xi^-,\xi_\perp)$，且
$p^\mu = (p^+\gamma, 0, p^-=M^2/(p^+\gamma)$；$M$ 是强子质量，$\gamma$ 是助推参数。
取 $\gamma\rightarrow\infty$ 即回到无限动量系。在这个新参考系中，Ioffe 时间为 $\nu = -p^+\xi^-$，
而 $-\xi^2= \xi_\perp^2$ 同样类空，二者都与 $\gamma$ 无关。此时 $h$ 成为近光锥关联函数
\begin{equation}\label{eq:lccorrel}
h(p^+\xi^-, \xi_\perp, \Lambda) =\langle p|\overline{\psi}(0,\xi^-,\xi_\perp)\Gamma L(0,\xi^-,\xi_\perp;0)\psi(0)|p\rangle,
\end{equation}
其中的规范连接现在同时具有沿光锥和横向方向的分量。洛伦兹对称性保证了 Eqs.~(\ref{eq:covitd})
与~(\ref{eq:lccorrel}) 的等同性，其中 $\xi_\perp = z$、$\xi^- = -zp_z/p^+$。相对 Ioffe 时间
$\nu$ 作傅里叶变换后直接得到光锥动量分数 $x$，因此如无限动量系中一贯的情形那样限于 $[-1,1]$。

于是可以从类空关联引入一个``赝''分布（pseudo-distribution），
\begin{equation}
{\cal P}(y, \xi_\perp, \Lambda) = \frac{1}{2\pi} \int^\infty_{-\infty}d\nu\ e^{-iy\nu}\ h(zp_z,z^2, \Lambda)\ ,
\end{equation}
其中 $\xi_\perp^2 = z^2$。在这一做法中，格点上的强子动量被重新解释为固定 $\xi_\perp$ 处的
Ioffe 时间。赝分布的物理在于：它是某一类横动量依赖部分子分布的傅里叶变换。由于 Eq.~(\ref{eq:lccorrel})
中夸克场之间有一个沿横向方向的规范连接，在 $A^+=0$ 规范（guage，原文如此）下分布中涉及无穷多个
横向极化的胶子。因此赝分布的部分子图像并不简单。特别地，动量分数 $y$ 不仅包含单个夸克的贡献，
还包含横向极化的物理胶子的贡献。只有当 $\xi_\perp=0$ 时 $y$ 才趋于夸克的光锥分数。

在这种对欧氏矩阵元的替代诠释中，把夸克完全放到光锥上——即夸克的横向坐标 $\xi_\perp = z$
趋于零——即可恢复物理夸克分布。在此极限下，为了保持有限的 Ioffe 时间 $\nu = zp_z$，又必须让
$p_z \rightarrow\infty$，与 LaMET 中的做法相同。因此，两种方法对获得精度部分子分布的要求
完全一致。

实践中必须使用小而非零的 $\xi_\perp$ 以避免光锥发散。于是可以建立一个与 LaMET 类似的因子化
定理（``小距离有效理论''），
\begin{equation}
{\cal P}(y, \xi_\perp,\tilde{\mu}) = \int \frac{dx}{x}
C(y/x, \xi_\perp\tilde{\mu},\tilde{\mu}/\mu) q(x, \mu) + {\cal O}(\xi^2_\perp),
\label{fac2}
\end{equation}
其中 $q$ 是物理部分子分布，$C$ 是依赖于具体格点正规化的匹配系数，$\tilde{\mu}$ 与 $\mu$
分别是 Ioffe 时间分布和物理 PDF 的重正化标度。要有连续极限，$\xi_\perp$ 应远大于格距 $a$，
但远小于 $1/\Lambda_{\text{QCD}}$。修正是按小 $\xi_\perp$ 展开的。在固定 $\xi_\perp$ 下通过上述
因子化公式提取部分子分布将会很有意义。

为了演示上述因子化定理，我们以同位旋矢量非极化夸克分布为例。在费曼规范下，我们在维数正规化
$D=4-2\epsilon$ 下单圈计算准 PDF 与物理 PDF 的夸克矩阵元。外夸克态取在壳且无质量，因此紫外
与共线发散分别由 $1/\epsilon_{UV}\ (\epsilon_{UV}>0)$ 与 $1/\epsilon_{IR}\ (\epsilon_{IR}<0)$
正规化。我们在坐标空间计算与文献~\cite{Xiong:2013bka} 相同的费曼图，得到
$h^{(1)}(zp^z,z^2,\mu,\epsilon)$（取 $\Gamma=\gamma^0$）与 $q^{(1)}(p^+z^-,\mu,\epsilon)$，
其中 $\mu$ 可视为重正化标度。费曼规则并非常规，因为每个不可约图中都有一个圈动量分量被作傅里叶
变换，例如，
\begin{align}
\mu^{2\epsilon}\int {d^d k\over (2\pi)^d} {1\over k^2(p-k)^2} e^{-i k^z z} &=  i\mu^{2\epsilon}\int_0^1 du \int {d^d k\over (2\pi)^d} {e^{-i k^z z}\over k_E^4} e^{ -i u p^z z}\ .
\end{align}
利用施温格参数化，可以解析求出上述积分，
\begin{align}
&i\mu^{2\epsilon}\int_0^1 du\int {d^d k\over (2\pi)^d} {e^{-i k^z z}\over k_E^4} e^{-i u p^z z}\nonumber\\
&= i\mu^{2\epsilon}\int_0^1 du\int_0^\infty d\alpha\ \alpha\int {d^d k\over (2\pi)^d} e^{-\alpha k_E^4-i k^z z} e^{-i u p^z z} \nonumber\\
&= {i\over 16\pi^2}(\pi z^2\mu^2)^{\epsilon_{IR}} \Gamma(-\epsilon_{IR})\int_0^1 du\  e^{-iu p^z z}\ .
\end{align}
$h^{(1)}(zp^z,z^2,\mu,\epsilon)$ 与 $q^{(1)}(p^+z^-,\mu,\epsilon)$ 的完整结果是
\begin{align}
&h^{(1)}(zp^z,z^2,\mu,\epsilon)\nonumber\\
=&{\alpha_sC_F\over 2\pi}\int_0^1 du \left[\left(1-u+\left({2u\over 1-u}\right)_+\right)\left(-\ln(z^2 \mu^2) - {1\over \epsilon'_{IR}}\right)\right.\nonumber\\
&\ \ \ \ \ \ \ \ + (1-u)\Big]e^{-iu p^z z}\nonumber\\
& + {\alpha_sC_F\over 2\pi}(izp^z)\int_0^1 du \int_0^1dt\ (2-u) (-\ln t^2) e^{-i(1-tu) p^z z} \nonumber\\
&  +\left[(1+\delta Z_\psi)+ {\alpha_sC_F\over 2\pi}\left(2\ln(z^2\mu^2) + {2\over \epsilon'_{UV}}+2\right)\right] e^{-i p^z z} \ ,
\end{align}

\begin{align}
&q^{(1)}(p^+z^-,\mu,\epsilon)\nonumber\\
=&{\alpha_sC_F\over 2\pi}\int_0^1 du \left(1-u+\left({2u\over 1-u}\right)_+\right)\left({1\over \epsilon'_{UV}}- {1\over \epsilon'_{IR}}\right)e^{iu p^+ z^-}\nonumber\\
& + (1+\delta Z_\psi)e^{i p^+ z^-} \ ,
\end{align}
其中 $Z_\psi = 1+ \delta Z_\psi$ 是夸克波函数重正化常数，并且我们作了替换
$1/\epsilon'_{UV,IR} = 1/\epsilon_{UV,IR}+\gamma_E+\ln\pi$。为保证矢量流守恒，$Z_\psi$ 在
在壳方案中给出为
\beq
Z_\psi = 1 - {\alpha_sC_F\over 2\pi}{1\over2}\left({1\over \epsilon'_{UV}}- {1\over \epsilon'_{IR}}\right) + O(\alpha_s^2)\ .
\eeq
然后把 Ioffe 时间 $p^z z$ 或 $p^+z^-$ 傅里叶变换到 $x$，得
\begin{align} \label{eq:quasipdf}
&{\cal P}^{(1)}(x,\xi_\perp,\mu,\epsilon)\nonumber\\
=&{\alpha_sC_F\over 2\pi}\left[\left({1+x^2\over 1-x}\right)_+\left(-\ln(\xi_\perp^2 \mu^2) - {1\over \epsilon'_{IR}}-1\right)  - \left(4\ln(1-x)\over 1-x\right)_+  \right.\nonumber\\
&+2(1-x)\Big]\theta(x)\theta(1-x)\nonumber\\
&  +\left[ 1+ {\alpha_sC_F\over 2\pi}\left({3\over2}\ln(\xi_\perp^2\mu^2) + {3\over2}{1\over \epsilon'_{UV}}+{3\over2}\right)\right]\delta(1-x) \ ,
\end{align}
而相应的光锥 PDF 是
\begin{align} \label{eq:lcpdf}
q^{(1)}(x,\mu,\epsilon) &= \delta(1-x) + {\alpha_s\over 2\pi}\left({1+x^2\over 1-x}\right)_+\left({1\over\epsilon'_{UV}} - {1\over \epsilon'_{IR}}\right)\ .
\end{align}
如预期那样，赝分布 ${\cal P}^{(1)}(x,\xi_\perp,\mu,\epsilon)$ 的支撑域限于 $0<x<1$。
若把 Eq.~(\ref{eq:quasipdf}) 与 Eq.~(\ref{eq:lcpdf}) 都在 $\overline{\text{MS}}$ 方案中重正化，
则可从 Eq.~(\ref{fac2}) 读出匹配系数
\begin{align}
&C\left(y,\xi_\perp\mu\right)\nonumber\\
=& \left[1+{\alpha_sC_F\over 2\pi}\left({3\over2}\ln(\xi_\perp^2\mu^2)+{3\over2}\right)\right] \delta(1-y)\nonumber\\
& + {\alpha_sC_F\over 2\pi} \left[-\left({1+y^2\over 1-y}\right)_+\left(\ln(\xi_\perp^2 \mu^2) +1 \right)- \left(4\ln(1-y)\over 1-y\right)_+ \right.\nonumber\\
&+2(1-y)\Big]\theta(y)\theta(1-y)\ .
\end{align}
\noindent 类似的因子化公式可以推广到格点正规化和非微扰重正化的情形。

文献~\cite{Radyushkin:2017cyf} 并未使用 Eq.~(\ref{fac2}) 的因子化公式从固定 $\xi_\perp$ 的格点
矩阵元中提取部分子分布，而是关注 Ioffe 时间分布可能的近似因子化性质，即
${\cal M}(\nu, z) \sim  g(\nu)h(z)$。若真如此，则 $g(\nu) \sim {\cal M}(\nu, z)/h(z)$ 在不同 $z$
下会呈现近似标度行为。这样便可从大 $z$ 计算得到小 $z$ 的 Ioffe 时间分布，无需非常大的动量来产生
大的 Ioffe 时间。的确，一次探索性的格点计算表明这种特殊的因子化或标度行为工作得相当好~\cite{Orginos:2017kos}。

这种标度或因子化行为固然有趣——正如深度非弹性散射中的部分子--强子对偶~\cite{Ji:1994br}
允许在小 $Q^2$ 下提取部分子分布一样——但它在实际的部分子分布精度计算中用途有限。探索性计算中
的标度行为只是在有限的运动学区域、以有限的精度被观察到的。部分子分布涉及各种大小的 Ioffe 时间，
特别是小 $x$ 分布涉及非常大的 Ioffe 时间，在那里标度更难检验，而且当前的结果肯定是不对的。
要得到正确的大 Ioffe 时间行为，需要在小的 $\xi_\perp$ 和更大的核子动量下计算，这与 LaMET 的要求
相同。此外，新的 $\xi_\perp$ 依赖的因子化在精度计算中必须加以修正。由于对这一因子化破坏的修正没有
系统的办法，这样的唯象方法不会成为 Eq.~(\ref{fac2}) 那类严格方法的替代品——后者中修正可以被量化。

当然，当从相同的格点矩阵元出发提取时，比较这两种看似不同的因子化方法（大动量 vs 小距离）给出的
部分子分布将是有意义的。重要的是以一致的方式量化两种方法的提取误差。我们强调，LaMET 是一个更为
普遍的框架，可以用费曼所提供的简单物理图像计算部分子物理。对于如何系统地、在可量化误差下计算所有
部分子可观测量——包括多部分子振幅与关联——它都有明确的方案。
